\section{Implementation}

This section describes how PESAD is realised: the tools adopted, the modular
architecture of the system, the representation of the knowledge base, the
mechanisms of the inference engine ---knowledge acquisition, uncertainty
management and explanation--- and the bilingual user interface.

\subsection{Tools Used}

PESAD is implemented in \emph{SWI-Prolog}, a mature and free implementation of
the Prolog logic-programming language. Prolog was chosen because the diagnostic
knowledge is naturally expressed as \emph{rules} (definite clauses) and because
its resolution mechanism provides backward chaining and backtracking out of the
box. The features that Prolog does not provide natively ---reasoning under
uncertainty and the explanation of the reasoning process--- are implemented
explicitly on top of the language, as described below.

\subsection{System Architecture}

The system is organised by \emph{responsibility} rather than as a single
monolithic file. This separation mirrors the conceptual architecture of an
expert system (knowledge base, inference engine, user interface) and makes each
concern independently readable and maintainable. Table~\ref{tab:modules} lists
the modules and the directives that load them are collected in
\texttt{start.pl}.

\begin{table}[h]
\centering
\caption{Modular organisation of PESAD.}
\label{tab:modules}
\small
\begin{tabular}{@{}lp{9.2cm}@{}}
\toprule
\textbf{Module} & \textbf{Responsibility} \\
\midrule
\texttt{start.pl} & Bootstrap: directives and loading of the modules. \\
\texttt{shell.pl} & User shell: REPL loop, language selection, command dispatch, presentation of results, and the \texttt{decode/2}\,/\,\texttt{explain/2} localization dispatch. \\
\texttt{ask.pl} & Knowledge acquisition: the \emph{criterion}, \emph{multiple} and \emph{selective} question protocols, with answer parsing and certainty elicitation. \\
\texttt{uncertainty.pl} & Certainty-factor calculus, certainty-based ordering and inference-level index bookkeeping. \\
\texttt{explanation.pl} & Printing of the proof tree (\emph{how} a conclusion was reached). \\
\texttt{memory.pl} & Working-memory management (facts and proof trees). \\
\texttt{utils.pl} & General-purpose helpers (output spacing, list operations). \\
\texttt{disorders.pl} & Knowledge base: codable Axis~I disorders (diagnostic goals) and entry points. \\
\texttt{criteria.pl} & Knowledge base: sub-disorders, types, symptomatic manifestations and features, and the questions. \\
\texttt{interface\_en.pl}, \texttt{interface\_it.pl} & English / Italian localization (\texttt{label/3}, \texttt{explanation/3}, \texttt{command/3}). \\
\bottomrule
\end{tabular}
\end{table}

\subsubsection{Knowledge Base}

The knowledge base is structured over \emph{four inference levels} that descend
from the abstract diagnoses to the concrete observations:

\begin{enumerate}[label=(\Roman*)]
    \item \textbf{Disorders} ---the codable Axis~I diagnoses, i.e. the goals of
    the system (\texttt{disorder/3});
    \item \textbf{Sub-disorders and types} ---non-codable building-block
    disorders and their specifiers (\texttt{disorder\_nc/4} and the
    \texttt{*\_type/4} predicates);
    \item \textbf{Symptomatic manifestations and features}
    (\texttt{symptomatic\_manifestation/4}, \texttt{symptomatic\_feature/4});
    \item \textbf{Elementary symptoms / facts} elicited from the user
    (\texttt{symptom/3}).
\end{enumerate}

Each rule carries, besides the patient and the resulting certainty factor, the
truth value it is asserted with. As an example, the rule for \emph{Panic
Disorder} requires the presence of the panic-disorder criteria and the
\emph{absence} of an associated substance-induced alteration and of an
associated general medical condition (the shared exclusion):

\begin{lstlisting}
disorder(panic_disorder, Patient, CF) :-
    disorder_nc(panic_disorder, Patient, CF1, true),
    exclude_substance_and_medical(Patient, CF2, CF3),
    certainty(panic_disorder, 1, [CF1,CF2,CF3], 0.97, CF, true).
\end{lstlisting}

The last argument of \texttt{certainty/6} is the truth value, the constant
\texttt{0.97} is the certainty of the rule itself, and the integer \texttt{1}
identifies the inference level, which is used by the explanation module to build
the proof tree.

\subsubsection{Inference Engine}

The inference engine is the part of the system that, given the facts stated by
the user, derives the diagnoses. It comprises three concerns: the acquisition of
information from the user, the management of uncertainty, and the explanation of
the reasoning.

\paragraph{Knowledge Acquisition.}
Questions are posed to the user through three protocols, implemented in
\texttt{ask.pl}.

\emph{Criterion ask} verifies a criterion made of one or more items, each a
yes/no question; a criterion of $N$ items may require that only $M$ of them
(with $1 \le M \le N$) be satisfied. Its certainty is obtained by taking the $M$
items with the highest certainty and then keeping the \emph{lowest} certainty
among them:
\begin{equation}
\cf_{\text{criterion}} \;=\; \min \big( \text{top-}M \{\, \cf_1, \dots, \cf_N \,\} \big).
\end{equation}
For instance, with $N=3$ items of certainties $\{0.4, 0.7, 0.6\}$ and $M=2$, the
two highest are $\{0.7, 0.6\}$ and the resulting certainty is $0.6$. If the user
answers \emph{no} to an item with certainty $c$, the item contributes the
complementary certainty $1-c$.

\emph{Multiple ask} lets the user choose one value from a list of options. If
the chosen answer is among those admitted by the question, the certainty equals
the one declared by the user; otherwise the residual certainty is distributed
over the remaining options:
\begin{equation}
\cf \;=\;
\begin{cases}
\cf_{\text{answer}} & \text{if the answer is admitted,}\\[4pt]
\dfrac{1-\cf_{\text{answer}}}{\,|\text{options}|-1\,} & \text{otherwise.}
\end{cases}
\end{equation}

\emph{Selective ask} is used for single, mutually exclusive choices (e.g. the
patient's age group): the admitted option succeeds with certainty $1.0$, while
any other choice makes the goal fail.

\paragraph{Uncertainty Management.}
Reasoning under uncertainty is implemented with \emph{certainty factors}
(\cf{}~$\in [0,1]$), managed in \texttt{uncertainty.pl}. The certainty of the
\emph{antecedent} of a rule, whose premises must all hold (a conjunction), is
the minimum of the premises' certainties:
\begin{equation}
\cf_{\text{ant}} \;=\; \min_i \cf_i .
\end{equation}
The certainty of the \emph{consequent} combines the antecedent certainty with
the certainty of the rule. For a goal asserted as true the two are multiplied;
for a negated goal the complement of the product is taken:
\begin{equation}
\cf_{\text{cons}} \;=\;
\begin{cases}
\cf_{\text{rule}} \cdot \cf_{\text{ant}} & \text{(positive goal),}\\[4pt]
1 - \cf_{\text{rule}} \cdot \cf_{\text{ant}} & \text{(negated goal).}
\end{cases}
\end{equation}
These two predicates are at the heart of the calculus:

\begin{lstlisting}
antecedent_certainty(Antecedent_CF_List, Antecedent_CF) :-
    min_list(Antecedent_CF_List, Antecedent_CF).

consequent_certainty(Goal, 1, Antecedent_CF, Rule_CF, Consequent_CF, true) :-
    Consequent_CF is Rule_CF * Antecedent_CF,
    ...
consequent_certainty(Goal, 1, Antecedent_CF, Rule_CF, Consequent_CF, false) :-
    Consequent_CF is 1 - (Rule_CF * Antecedent_CF),
    ...
\end{lstlisting}

While computing the consequent, each rule is asserted (as \texttt{rule\_I/4},
\texttt{rule\_II/5} or \texttt{rule\_III/6}, according to its inference level)
together with a hierarchical index. These assertions form the proof tree that
the explanation module later traverses.

\paragraph{Certainty factors versus fuzzy logic.}
The conjunction of the premises is, in fact, a \emph{fuzzy} operation: PESAD
evaluates it with a configurable t-norm, selected at start-up and stored in the
dynamic fact \texttt{uncertainty\_method/1}. Two methods are provided:
\begin{itemize}
  \item \texttt{cf} (the default) uses the \emph{minimum} (G\"odel t-norm),
  i.e. the certainty of a criterion equals that of its \emph{weakest} satisfied
  premise: $\cf_{\text{ant}} = \min_i \cf_i$;
  \item \texttt{fuzzy} uses the \emph{product} t-norm, i.e. the algebraic
  product of the premises: $\cf_{\text{ant}} = \prod_i \cf_i$.
\end{itemize}
The remaining steps (rule chaining $\cf_{\text{rule}}\cdot\cf_{\text{ant}}$ and
the negation $1-x$) are shared. The two methods are compared on a suite of
clinical vignettes in Chapter~\ref{chap:validation}; both rank the correct
diagnosis first, but the product t-norm systematically \emph{under-rates
polythetic disorders}: a fully met post-traumatic stress disorder, which
combines seven symptom clusters,
collapses to $28\%$ because every premise multiplies the result, whereas the
minimum keeps the confidence at a clinically plausible $66\%$ regardless of how
many criteria a disorder has. For a diagnostic task in which a syndrome is ``as
strong as its weakest satisfied criterion'', the minimum is the more faithful
operator, and it is therefore kept as the default; the product t-norm remains
available for comparison and experimentation.

A richer fuzzy treatment is also provided as a third method
(\texttt{fuzzy\_membership}, in \texttt{fuzzy.pl}): for the graded clinical
variables ---symptom duration, panic-attack frequency and degree of insight---
the clinician enters a \emph{crisp} value (e.g. the duration in months) instead
of choosing a discrete option, and a trapezoidal membership grade becomes the
certainty. This replaces the hard diagnostic cut-offs with smooth boundaries:
for instance, against the ``at least six months'' criterion a five-month
duration scores $0.67$ rather than abruptly failing, while the conjunction is
still aggregated with the minimum. Table~\ref{tab:duration-membership}
illustrates the membership of the ``at least six months'' fuzzy set.

\begin{table}[h]
\centering
\caption{Membership grade of the ``duration at least 6 months'' fuzzy set.}
\label{tab:duration-membership}
\small
\begin{tabular}{@{}lcccccc@{}}
\toprule
\textbf{duration (months)} & 2 & 3 & 4 & 5 & 6 & 8 \\
\midrule
\textbf{membership} & 0\% & 0\% & 33\% & 67\% & 100\% & 100\% \\
\bottomrule
\end{tabular}
\end{table}

\paragraph{Explanation Module.}
A good expert system must justify its behaviour. PESAD does so in two ways,
implemented in \texttt{ask.pl} (the \emph{why}) and \texttt{explanation.pl} (the
\emph{how}).

\emph{Why} questions: while answering a yes/no or multiple-choice question, the
user may reply \texttt{why} (\texttt{perche}) to receive an explanation of the
diagnostic reason behind the question; the same question is then asked again.

\emph{How} questions: after a diagnosis the user may ask how it was reached. The
predicate \texttt{print\_tree/1} walks the asserted rules across the three
inference levels and prints, indented by level, each (sub-)conclusion with its
certainty, distinguishing what was \emph{deduced} (intermediate rules) from what
was \emph{stated} by the user (leaf facts).

\subsubsection{User Interface}

The user interacts with the system through a textual shell (\texttt{shell.pl})
that implements a read--evaluate--print loop. At start-up the user selects the
interface language; the dynamic fact \texttt{language/1} then drives every
message through the localization dispatch:

\begin{lstlisting}
decode(Key, Text) :- language(Language), label(Language, Key, Text), !.
decode(Key, Key).
\end{lstlisting}

All user-visible text is therefore stored as language-neutral \emph{keys} that
each locale file (\texttt{interface\_en.pl}, \texttt{interface\_it.pl}) maps to a
\texttt{label/3} string and, where relevant, to an \texttt{explanation/3} string.
The command vocabulary is localized in the same way through \texttt{command/3},
so that the user may type either the English or the Italian command words while
the engine works on a single set of internal command atoms
(\texttt{investigation}, \texttt{control}, \texttt{facts}, \texttt{help},
\texttt{clean}, \texttt{quit}). Adding a new language amounts to providing a new
locale file, with no change to the engine or the knowledge base.
